\begin{tikzpicture}[x=1cm, y=1cm, font=\scriptsize]

  % Keep the graphic's bounding box symmetric around the origin for stable centering.
  \path[use as bounding box] (-7.0,-7.0) rectangle (7.0,7.0);

  % ========== STYLE DEFINITIONS ==========
  \tikzset{
    postit inner/.style={
      rectangle, rounded corners=1pt, draw=black, line width=0.4pt,
      fill=red!10, minimum width=18mm, minimum height=12mm,
      align=center, text=mmInk, font=\scriptsize\itshape, inner sep=1.5mm
    },
    postit middle/.style={
      rectangle, rounded corners=1pt, draw=black, line width=0.4pt,
      fill=orange!15, minimum width=18mm, minimum height=12mm,
      align=center, text=mmInk, font=\scriptsize\itshape, inner sep=1.5mm
    },
    postit outer/.style={
      rectangle, rounded corners=1pt, draw=black, line width=0.4pt,
      fill=green!10, minimum width=18mm, minimum height=12mm,
      align=center, text=mmInk, font=\scriptsize\itshape, inner sep=1.5mm
    },
  }

  % ========== THREE CONCENTRIC CIRCLES (filled outside-in for ring effect) ==========
  \fill[green!10] (0,0) circle (6.8cm);
  \draw[draw=black, line width=0.8pt] (0,0) circle (6.8cm);
  \fill[orange!15] (0,0) circle (5.0cm);
  \draw[draw=black, line width=0.8pt] (0,0) circle (5.0cm);
  \fill[red!10] (0,0) circle (2.6cm);
  \draw[draw=black, line width=0.8pt] (0,0) circle (2.6cm);

  % ========== CENTER: Person icon ==========
  % Head
  \fill[blue!40!black] (0,0.45) circle (0.35cm);
  % Body
  \fill[blue!40!black] (-0.5,-0.15) to[out=80,in=180] (0,0.15)
    to[out=0,in=100] (0.5,-0.15)
    to[out=-30,in=-150] (0.5,-0.7) -- (-0.5,-0.7)
    to[out=-30,in=210] (-0.5,-0.15);

  % ========== INNER CIRCLE: Primäre Nutzende ==========
  \node[postit inner, minimum width=22mm] at (0,1.65) {Nutzer: vitale\\Senioren\\(Ausflügler:in)};
  \node[postit inner, minimum width=22mm] at (0,-1.65) {soziales\\Umfeld};

  % ========== MIDDLE CIRCLE: Schlüssel-Partner ==========
  % Top area
  \node[postit middle] at (-2.0, 3.85) {bereits öV\\fahrende\\Senioren};
  \node[postit middle] at (2.0, 3.85) {Bund \&\\Kantone};

  % Left area
  \node[postit middle, minimum width=26mm, minimum height=22mm, font=\tiny\itshape] at (-3.5, 2.3)
    {``Invisible Influencers''\\Ärzte (Führerschein-\\Thema), Gemeinden\\(Infrastruktur,\\Angebote),\\Versicherungen,\\Medien / gesellsch.\\Bild};

  % Right area
  \node[postit middle, minimum width=20mm, font=\normalsize\itshape] at (3.85, 2.55) {SBB};
  \node[postit middle, minimum width=20mm, font=\normalsize\itshape] at (3.85, 0.18) {ZSG};

  % Mid-left
  \node[postit middle] at (-2.95, -0.45) {Tourismus-\\Vereinigungen};
  \node[postit middle] at (-2.95, -2.55) {Bahnhöfe\\(Shops,\\Restaurants, ...)};

  % Bottom area
  \node[postit middle, minimum width=22mm] at (-1.35, -3.65) {Bewerter:\\CAS Profs.\\(Sandro \& Nina)};
  \node[postit middle] at (1.35, -3.65) {weitere ZVV\\Dienstleister};
  \node[postit middle] at (3.2, -3.2) {Weitere\\Verkehrs-\\verbünde};
  \node[postit middle] at (3.65, -1.35) {Auftraggeber:\\ZVV (Pascal)};

  % ========== OUTER CIRCLE / OUTSIDE: Erweitertes Umfeld ==========
  \node[postit outer] at (0, 6.2) {Pendler\\$< 65$};
  \node[postit outer] at (4.6, 5.4) {Gesellschaft\\Schweiz};
  \node[postit outer] at (-5.7, 0) {Autolobby};
  \node[postit outer] at (5.7, 0) {Umwelt-\\Bewusste};
  \node[postit outer] at (0, -6.2) {Freizeitanbieter};

\end{tikzpicture}
